
\chapter{Fourier 级数逼近与黎曼对数余弦逼近的比较}\label{ch:fourier_vs_riemann}

\section{引言}\label{sec:fourier_vs_riemann_intro}

上一章的显式公式将 $\psi(x)$（跳跃点取左右平均后为 $\psi_0$）表示为主项与零点级数
$\sum_{\rho}x^{\rho}/\rho$ 之和。零点共轭成对时，该级数可合并为以 $\gamma=\operatorname{Im}\rho$
为对数尺度频率的实余弦型振荡。

本章只做一件事：把上述振荡写成 $\cos(\gamma\log x-\arg\rho)$ 的语言，并与周期 Fourier 逼近
$\cos(kx+\phi)$ 对照，说明二者原理不同。不给出新的素数定理。

\begin{definition}{$\pi_0(x)$ 的正则化}{}\label{def:pi0}
\[
\pi_0(x)=\frac{\pi(x^-)+\pi(x^+)}{2}.
\]
非素数处 $\pi_0=\pi$；素数 $p$ 处 $\pi_0(p)=\pi(p)-\frac12$。
显式公式在跳跃点收敛到 $\pi_0$，与 Fourier 级数在间断点收敛到左右均值同型。
\end{definition}

平凡零点项 $-\frac12\log(1-x^{-2})$ 经 $\psi\to\pi$ 后为 $O(1/\log x)$，本章近似中可略；
精确计算见第~\ref{ch:riemann_explicit_formula}~章。

% ============================================================
\section{零点对的余弦合并}
% ============================================================

设 $\rho=\beta+i\gamma$（$\gamma>0$）。由 $\zeta$ 的共轭对称，$\overline{\rho}=\beta-i\gamma$ 亦为零点
（函数方程给出 $\rho\leftrightarrow 1-\rho$，与共轭是不同对称）。$|\overline{\rho}|=|\rho|$，故
\begin{align*}
\frac{x^{\rho}}{\rho}+\frac{x^{\overline{\rho}}}{\overline{\rho}}
&= \frac{x^{\beta}}{|\rho|}
   \bigl(e^{i(\gamma\log x-\arg\rho)}+e^{-i(\gamma\log x-\arg\rho)}\bigr)
= \frac{2x^{\beta}}{|\rho|}\cos(\gamma\log x-\arg\rho),
\end{align*}
其中 $\arg\rho=\arctan(\gamma/\beta)\in(0,\pi/2)$（$\beta,\gamma>0$）。即
\begin{equation}
\boxed{\frac{x^{\rho}}{\rho}+\frac{x^{\overline{\rho}}}{\overline{\rho}}
=\frac{2x^{\beta}}{|\rho|}\cos(\gamma\log x-\arg\rho)}
\label{eq:cosine_pair}
\end{equation}
$\gamma$ 为对数尺度上的频率；$\beta$ 控制振幅 $x^{\beta}$；RH 下 $\beta=1/2$，
$\arg\rho=\arctan(2\gamma)$（$\gamma\gg 1/2$ 时接近 $\pi/2$）。

若 RH 成立，按 $\gamma>0$ 配对求和得
\begin{equation}
\boxed{\psi_0(x)=x-\sum_{\gamma>0}\frac{2x^{1/2}}{|\rho|}\cos(\gamma\log x-\arg\rho)
-\log(2\pi)-\frac12\log(1-x^{-2})}
\label{eq:cosine_psi}
\end{equation}
（$x>1$，对称截断；非素数幂处 $\psi_0=\psi$）。
每一项在 $\log x$ 上周期为 $2\pi/\gamma$；诸 $\gamma$ 一般非等距，故不是单一基频的谐波族。

% ============================================================
\section{对数余弦与标准余弦}
% ============================================================

\begin{center}
\renewcommand{\arraystretch}{1.45}
\begin{tabular}{@{}lll@{}}
\toprule
 & 标准余弦 $A\cos(kx+\omega)$ & 对数余弦 $A\cos(\gamma\log x+\omega)$ \\
\midrule
相位自变量 & $x$ & $\log x$ \\
波峰间距 & 等差：$\Delta x=2\pi/k$ & 等比：$x_{n+1}/x_n=e^{2\pi/\gamma}$ \\
\bottomrule
\end{tabular}
\end{center}

例：$\gamma_1\approx 14.1347$（RH），公比 $e^{2\pi/\gamma_1}\approx 1.560$，
即 $x$ 每增大约 $56\%$ 完成一周振荡。令 $\tau=\log x$，则
$\cos(\gamma\log x+\omega)=\cos(\gamma\tau+\omega)$ 在 $\tau$ 轴上为常频余弦。

由 $\psi\to J\to\pi$（第~\ref{ch:riemann_explicit_formula}~章与附录~\ref{app:abel_stieltjes}），
在 RH、对称截断并略去 $n\ge 2$ 的 M\"obius 项与平凡零点后，
\begin{equation}
\pi_0(x)-\Li(x)
\approx -\sum_{k=1}^{N}\frac{2\sqrt{x}}{|\rho_k|\log x}
\cos(\gamma_k\log x-\arg\rho_k),\qquad N\to\infty.
\label{eq:riemann}
\end{equation}
振幅 $2\sqrt{x}/(|\rho_k|\log x)$ 缓变；频率 $\{\gamma_k\}$ 一般不是某一基频的整数倍，
故上式是非等距余弦叠加，不是标准周期 Fourier 级数。

% ============================================================
\section{对数余弦的有限截断：\texorpdfstring{$N=3$}{N=3}}\label{sec:n3_principle}
% ============================================================

记
\begin{equation}
R_k(x)=\frac{2\sqrt{x}}{|\rho_k|\log x}\cos(\gamma_k\log x-\arg\rho_k),
\quad k=1,2,3,\ldots
\label{eq:Rk}
\end{equation}
（RH 下 $\beta=1/2$）。则式~\eqref{eq:riemann} 的有限截断为
$\pi_0(x)\approx\Li(x)-\sum_{k=1}^{N}R_k(x)$，仅含非平凡零点主项，
略去平凡零点及 M\"obius 反演中 $n\ge 2$ 的项，适用于定性展示而非精确计算。
单看一项只能说明某一 $\gamma$ 的修正；取前几个不同 $\gamma$ 叠加，
才能看出非等距频率如何共同修正主项。本节取 $N=3$ 作有限截断示意
（对照图见~\ref{fig:n3-logcos}）；系统加大 $N$ 见第~\ref{ch:riemann_numerical_approximation}~章。

\subsection{主项基准：\texorpdfstring{$\Li_2(x)$}{Li2(x)}}

采用 $\Li(x)$ 的渐近展开前两项作为逼近 $\pi_0(x)$ 的平滑基准。由
$\operatorname{li}(x)$ 的渐近级数
$\operatorname{li}(x)\sim\sum_{k=0}^{\infty}\frac{k!\,x}{(\log x)^{k+1}}$，
取前两项并利用 $\Li(x)=\operatorname{li}(x)-\operatorname{li}(2)$，得
\[
\Li_2(x)=\frac{x}{\log x}+\frac{x}{(\log x)^2}-\operatorname{li}(2),
\]
其中 $\operatorname{li}(2)=\int_0^2\frac{dt}{\log t}\approx 1.045$ 是固定常数
（注意 $\operatorname{li}(2)\neq\log 2\approx 0.693$；部分简化处理中将此项写为 $-\log 2$
系近似处理，误差约 $0.35$，在 $x\in[2,50]$ 范围内对图形影响较小，但在精确计算
中应使用 $\operatorname{li}(2)\approx 1.045$）。
截断误差为 $O\!\bigl(x/(\log x)^3\bigr)$，对 $x\le 50$ 而言绝对误差小于 $0.3$。

为使数值图形与既有作图约定一致，下面的 pgfplots 代码仍采用近似值
$-\log 2\approx -0.693$（最终影响约半个单位的垂直偏移）。读者如需精确计算，
应将该常数替换为 $-\operatorname{li}(2)\approx -1.045$。

图示范围取 $x\ge 2$，是因为振幅因子 $1/\log x$ 在 $x\to 1$ 时趋于无穷
（$\log x\to 0$），余弦修正量 $R_k(x)$ 在 $x\in(1,2)$ 附近急剧发散，不具有物理意义；
而 $\pi(x)=0$（$x<2$），显式公式在该区间内的近似意义也有限。故本节图示均从 $x=2$ 开始。

\subsection{\texorpdfstring{$N=3$}{N=3}：三零点叠加}

下图上方：三个 $R_k$ 及其合成线；下方：右连续 $\pi(x)$（示意）与 $N=3$ 逼近曲线。
图中坐标在素数处取 $\pi(p)$，而非半整数 $\pi_0(p)=\pi(p)-\frac12$；
理论目标函数是跳跃点均值 $\pi_0$，与 $\pi$ 仅在素数处相差 $\frac12$，宏观图几乎重合。

\paragraph{零点数据来源}
本数值模拟采用的前三个非平凡零点数据：
\[
\gamma_1\approx 14.134725,\quad
\gamma_2\approx 21.022040,\quad
\gamma_3\approx 25.010858,
\]
相应模值和辐角已按 $|\rho_k|=\sqrt{(1/2)^2+\gamma_k^2}$ 和
$\arg\rho_k=\arctan(2\gamma_k)$（RH 条件下 $\beta=1/2$）计算。
零点数据可参见 Odlyzko 的零点表\footnote{%
  A. M. Odlyzko, ``Tables of zeros of the Riemann $\zeta$ function,''
  \url{https://www.dtc.umn.edu/~odlyzko/zeta_tables/}}
或 LMFDB (L-Functions and Modular Forms Database)。

\begin{figure}[h!]
\centering
\begin{tikzpicture}
  \pgfmathsetmacro{\gA}{14.134725}\pgfmathsetmacro{\absA}{14.143}\pgfmathsetmacro{\argA}{1.534}
  \pgfmathsetmacro{\gB}{21.022040}\pgfmathsetmacro{\absB}{21.028}\pgfmathsetmacro{\argB}{1.547}
  \pgfmathsetmacro{\gC}{25.010858}\pgfmathsetmacro{\absC}{25.016}\pgfmathsetmacro{\argC}{1.551}
  \pgfmathsetmacro{\lnTwo}{0.693147}
  \begin{groupplot}[
    group style={group size=1 by 2, vertical sep=3.5cm,
                 x descriptions at=edge bottom},
    width=14cm, height=7cm, xmin=2, xmax=50,
    axis lines=middle, grid=both, grid style={gray!10}, xlabel={$x$}]
    %--- 上图：三个余弦修正量 ---
    \nextgroupplot[ymin=-1.5, ymax=1.5, ylabel={$R_k(x)$},
      legend pos=south west,
      legend style={font=\scriptsize,fill=white,yshift=-0.5cm,xshift=0.6cm},
      title={零点振荡修正量 ($R_k$) 与净合成线 ($\sum_3 R_k$)}]
    \addplot[brown!20!black,opacity=0.4,domain=2:50,samples=600]
      {(2*sqrt(x)/(\absA*ln(x)))*cos(deg(\gA*ln(x)-\argA))};
    \addlegendentry{$R_1$}
    \addplot[blue,opacity=0.4,domain=2:50,samples=600]
      {(2*sqrt(x)/(\absB*ln(x)))*cos(deg(\gB*ln(x)-\argB))};
    \addlegendentry{$R_2$}
    \addplot[green!60!black,opacity=0.4,domain=2:50,samples=600]
      {(2*sqrt(x)/(\absC*ln(x)))*cos(deg(\gC*ln(x)-\argC))};
    \addlegendentry{$R_3$}
    \addplot[red,thick,domain=2:50,samples=1000] {
      (2*sqrt(x)/(\absA*ln(x)))*cos(deg(\gA*ln(x)-\argA)) +
      (2*sqrt(x)/(\absB*ln(x)))*cos(deg(\gB*ln(x)-\argB)) +
      (2*sqrt(x)/(\absC*ln(x)))*cos(deg(\gC*ln(x)-\argC))};
    \addlegendentry{$R_{\Sigma_3}=\sum_{k=1}^3 R_k$}
    %--- 下图：红色合成与蓝色逼近 ---
    \nextgroupplot[ymin=-3, ymax=18, ylabel={$\pi(x)$},
      legend pos=north west,
      legend style={font=\scriptsize,fill=white,xshift=0.6cm},
      title={红色：$R_{\Sigma_3}$（零点合成）与蓝色：$\Li_2-R_{\Sigma_3}$（$N=3$ 逼近）对比}]
    \addplot[const plot,black,thick] coordinates {
      (2,1)(3,2)(5,3)(7,4)(11,5)(13,6)(17,7)(19,8)
      (23,9)(29,10)(31,11)(37,12)(41,13)(43,14)(47,15)(50,15)};
    \addlegendentry{右连续 $\pi(x)$（示意）}
    \addplot[green!50!black,domain=2:50,samples=150]
      {x/ln(x)+x/(ln(x)^2)-\lnTwo};
    \addlegendentry{主项 $\Li_2(x)$}
    \addplot[red,thick,domain=2:50,samples=800] {
      (2*sqrt(x)/(\absA*ln(x)))*cos(deg(\gA*ln(x)-\argA)) +
      (2*sqrt(x)/(\absB*ln(x)))*cos(deg(\gB*ln(x)-\argB)) +
      (2*sqrt(x)/(\absC*ln(x)))*cos(deg(\gC*ln(x)-\argC))};
    \addlegendentry{$R_{\Sigma_3}$（红色，振幅较大）}
    \addplot[blue,thick,domain=2:50,samples=800] {
      (x/ln(x)+x/(ln(x)^2)-\lnTwo) - (
        (2*sqrt(x)/(\absA*ln(x)))*cos(deg(\gA*ln(x)-\argA)) +
        (2*sqrt(x)/(\absB*ln(x)))*cos(deg(\gB*ln(x)-\argB)) +
        (2*sqrt(x)/(\absC*ln(x)))*cos(deg(\gC*ln(x)-\argC)))};
    \addlegendentry{$\Li_2-R_{\Sigma_3}$（蓝色，振荡较小）}
  \end{groupplot}
\end{tikzpicture}
\caption{频谱合成对素数阶梯的几何逼近（$N=3$，$x\ge 2$）。
上图：三个不同频率的对数余弦修正量 $R_k$ 与净合成线
$R_{\Sigma_3}=\sum_{k=1}^{3}R_k$；
下图：右连续 $\pi(x)$（示意）、主项 $\Li_2$、合成修正 $R_{\Sigma_3}$
与 $N=3$ 逼近 $\Li_2-R_{\Sigma_3}$。
理论目标在跳跃点为 $\pi_0$；略去平凡零点与 $n\ge 2$ 的 M\"obius 项。
图中 $\Li_2$ 常数项取 $-\log 2$（见正文）。
\textbf{注意：}为清楚显示三零点振荡合成的细节，上下两图垂直轴高度和水平轴长度的比例不同。}
\label{fig:n3-logcos}
\end{figure}

观察可知，上图中红色曲线 $R_{\Sigma_3}$ 的振幅（约 $\pm 0.6$）相对于纵轴范围
$[-1.5,1.5]$ 显得较大；而在下图中，由于纵轴范围扩大至约 $[-3,18]$，
同一条红色曲线的振幅在视觉上被压缩成一条几乎平坦的细线。
这纯粹是纵轴缩放比例不同造成的视觉差异，而非数学上的振幅改变。

观察可知：随着参与合成的零点数量 $N$ 增加，逼近曲线越来越贴近素数阶梯
（图为右连续 $\pi(x)$；精确公式在跳跃点收敛到 $\pi_0$）。
上方图中三个不同频率的 $R_k$ 叠加产生净合成线 $R_{\Sigma_3}$，
正是这条合成线通过修正主项，逐步逼近素数的阶梯位置。

以上仅为 $N=3$ 的有限截断示意。第~\ref{ch:riemann_numerical_approximation}~章将逐步增加
零点个数逼近 $\pi_0(x)$。下文先给出对照所需的 Fourier 级数与 Gibbs 要点
（第~\ref{ch:integral_transforms}~章是 Fourier \textbf{变换}，对象不同），
再以 $N=100$ 并置对照（图~\ref{fig:fourier-approx}、\ref{fig:riemann-logcos-approx}）。

% ============================================================
\section{Fourier 级数与 Gibbs 现象}\label{sec:fourier_gibbs}
% ============================================================

第~\ref{ch:integral_transforms}~章的 Fourier \textbf{变换}是 $\R$ 上的积分；
本章对照用的是周期函数的 \textbf{Fourier 级数}（等距三角谐波叠加）。二者相关，但本章只需级数一侧的标准事实。

\begin{definition}{周期函数的 Fourier 级数}{}\label{def:fourier_series}
设 $f$ 在 $(-L,L)$ 上可积，并按周期 $2L$ 延拓。Fourier 系数为
\begin{align*}
a_0&=\frac{1}{L}\int_{-L}^{L}f(t)\,\mathrm{d}t,\qquad
a_n=\frac{1}{L}\int_{-L}^{L}f(t)\cos\frac{n\pi t}{L}\,\mathrm{d}t,\\
b_n&=\frac{1}{L}\int_{-L}^{L}f(t)\sin\frac{n\pi t}{L}\,\mathrm{d}t
\qquad (n\ge 1).
\end{align*}
形式级数
\begin{equation}
f(x)\sim\frac{a_0}{2}+\sum_{n=1}^{\infty}
\Bigl(a_n\cos\frac{n\pi x}{L}+b_n\sin\frac{n\pi x}{L}\Bigr)
\label{eq:fourier}
\end{equation}
称为 $f$ 的 Fourier 级数；频率为基频 $\pi/L$ 的整数倍。
前 $N$ 项之和（含 $a_0/2$）
\begin{equation}
S_N(x)=\frac{a_0}{2}+\sum_{n=1}^{N}
\Bigl(a_n\cos\frac{n\pi x}{L}+b_n\sin\frac{n\pi x}{L}\Bigr)
\label{eq:fourier_partial}
\end{equation}
称为 \textbf{Fourier 部分和}：先定周期与等距谐波，再由 $f$ 投影得系数。
\end{definition}

若 $f$ 分段光滑，则在连续点通常 $S_N\to f$；在跳跃点 $x_0$，
$S_N(x_0)$ 收敛到左右极限的算术平均
$\bigl(f(x_0^-)+f(x_0^+)\bigr)/2$。
这与定义~\ref{def:pi0} 中 $\pi_0$ 取左右均值同型。

\begin{definition}{Gibbs 现象}{}\label{def:gibbs}
用 $S_N$ 逼近带跳跃的周期函数时，跳跃附近出现稳定的 \textbf{过冲}与\textbf{欠冲}：
过冲幅度约为跳跃高度的 $9\%$ 量级，随 $N\to\infty$ \textbf{不趋于零}；
波纹宽度随 $N$ 变窄。跳跃点本身仍收敛到左右均值。
该现象（Wilbraham--Gibbs）来自部分和与 \textbf{Dirichlet 核} 的卷积结构。
\end{definition}

有限个零点截断显式公式时，素数幂附近也可有起伏，同属「截断级数逼近阶梯」；
但截断核与 Dirichlet 核不同，过冲比例\textbf{不必}等于上述 $9\%$，不宜混称为同一定理。
下一节 $N=100$ 的 Fourier 图中可见经典过冲；对数余弦图中的起伏应归因于零点截断。

% ============================================================
\section{图示对照：\texorpdfstring{$N=100$}{N=100}}\label{subsec:comparison}
% ============================================================

前面以 $N=3$ 说明对数余弦的有限截断，并以 Fourier 部分和~\eqref{eq:fourier_partial}
与定义~\ref{def:gibbs} 作为读图工具。
本节在同一窗口 $x\in[10,35]$、$y\in[-1,6]$ 上
把\textbf{周期 Fourier 级数} 与\textbf{黎曼对数余弦} 各取 $N=100$ 并置：红线均为 $\pi(x)-\pi(16)$；
蓝线分别为 $S_{100}$ 与前 $100$ 个非平凡零点的重构，且仅画 $17\le x\le 31$
（第 $1$ 至第 $5$ 台阶）。此处 $N$ 的含义不同：谐波个数对零点个数
（图中 Fourier 侧周期取 $L=30$，$t=x-10$）。

\begin{figure}[h!]
\centering
\begin{tikzpicture}
\begin{axis}[
    width=14.0cm, height=9cm,
    xmin=10, xmax=35,
    ymin=-1.0, ymax=6.0,
    xlabel={$x$},
    ylabel={},
    title={\large\bf Fourier 级数（$N=100$）逼近 $\pi(x)-\pi(16)$},
    xtick={10,15,17,19,23,25,29,31,35},
    xticklabels={$10$,$15$,$17$,$19$,$23$,$25$,$29$,$31$,$35$},
    ytick={0,1,2,3,4,5},
    yticklabels={$0$,$1$,$2$,$3$,$4$,$5$},
    grid=both,
    grid style={line width=0.25pt, draw=gray!25},
    major grid style={line width=0.4pt, draw=gray!45},
    axis lines=left,
    tick align=outside,
    legend pos=north west,
    legend style={font=\small, fill=white, fill opacity=0.9, draw opacity=1},
    tick label style={font=\small},
    clip=true,
]
\addplot[red!80!black, very thick, solid, const plot]
    coordinates {(10,0)(17,1)(19,2)(23,3)(29,4)(31,5)(35,5)};
\addlegendentry{真实阶梯函数 $\pi(x)-\pi(16)$}
\addplot[blue!80!black, thick, line width=1.2pt,
    restrict x to domain=17:31]
    table[x index=0, y index=1, col sep=space]
    {fourier_data.dat};
\addlegendentry{Fourier 部分和 $S_{100}(x)$}
\draw[gray!50, very thin, dashed] (axis cs:17, -1.0) -- (axis cs:17, 6.0);
\draw[gray!50, very thin, dashed] (axis cs:19, -1.0) -- (axis cs:19, 6.0);
\draw[gray!50, very thin, dashed] (axis cs:23, -1.0) -- (axis cs:23, 6.0);
\draw[gray!50, very thin, dashed] (axis cs:29, -1.0) -- (axis cs:29, 6.0);
\draw[gray!50, very thin, dashed] (axis cs:31, -1.0) -- (axis cs:31, 6.0);
\node[font=\small, align=center, text=blue!70!black,
      fill=white, fill opacity=0.88, inner sep=2pt]
    at (axis cs:28, 0.0)
    {$S_N(x)=c_0+\sum_{n=1}^{N}\!\left[a_n\cos\!\tfrac{2\pi n t}{L}+b_n\sin\!\tfrac{2\pi n t}{L}\right]$};
\end{axis}
\end{tikzpicture}
\caption{周期 $L=30$、$N=100$ 的 Fourier 部分和逼近 $\pi(x)-\pi(16)$（$t=x-10$）。
跳跃处可见约 $9\%$ 量级过冲（Gibbs 现象）。}
\label{fig:fourier-approx}
\end{figure}

\clearpage

\begin{figure}[h!]
\centering
\begin{tikzpicture}
\begin{axis}[
    width=14.0cm, height=9cm,
    xmin=10, xmax=35,
    ymin=-1.0, ymax=6.0,
    xlabel={$x$},
    ylabel={},
    title={\large\bf 黎曼对数余弦（$N=100$）逼近 $\pi(x)-\pi(16)$},
    xtick={10,15,17,19,23,25,29,31,35},
    xticklabels={$10$,$15$,$17$,$19$,$23$,$25$,$29$,$31$,$35$},
    ytick={0,1,2,3,4,5},
    yticklabels={$0$,$1$,$2$,$3$,$4$,$5$},
    grid=both,
    grid style={line width=0.25pt, draw=gray!25},
    major grid style={line width=0.4pt, draw=gray!45},
    axis lines=left,
    tick align=outside,
    legend pos=north west,
    legend style={font=\small, fill=white, fill opacity=0.9, draw opacity=1},
    tick label style={font=\small},
    clip=true,
]
\addplot[red!80!black, very thick, solid, const plot]
    coordinates {(10,0)(17,1)(19,2)(23,3)(29,4)(31,5)(35,5)};
\addlegendentry{真实阶梯函数 $\pi(x)-\pi(16)$}
\addplot[blue!70!green!80!black, thick, line width=1.2pt,
    restrict x to domain=17:31, restrict y to domain=0:100]
    table[x index=0, y index=1, col sep=space]
    {riemann_data.dat};
\addlegendentry{显式公式 $\widehat{\pi}_{100}(x)$}
\draw[gray!50, very thin, dashed] (axis cs:17, -1.0) -- (axis cs:17, 6.0);
\draw[gray!50, very thin, dashed] (axis cs:19, -1.0) -- (axis cs:19, 6.0);
\draw[gray!50, very thin, dashed] (axis cs:23, -1.0) -- (axis cs:23, 6.0);
\draw[gray!50, very thin, dashed] (axis cs:29, -1.0) -- (axis cs:29, 6.0);
\draw[gray!50, very thin, dashed] (axis cs:31, -1.0) -- (axis cs:31, 6.0);
\node[font=\small, align=center, text=blue!70!green!80!black,
      fill=white, fill opacity=0.88, inner sep=2pt]
    at (axis cs:28, 0.0)
    {$\widehat{\pi}_N(x)=\Li(x)-C$\\[3pt]
     $
      {}-\frac{2\sqrt{x}}{\log x}
      \sum_{n=1}^{N}
      \frac{\cos(\gamma_n\log x-\arg\rho_n)}{|\rho_n|}$};
\end{axis}
\end{tikzpicture}
\caption{黎曼对数余弦逼近（前 $N=100$ 个非平凡零点）：蓝线 $\widehat{\pi}_N$ 与
式~\eqref{eq:riemann} 等价（一写误差、一写重构）；红线 $\pi(x)-\pi(16)$。
$C$ 仅使 $\widehat{\pi}_N(16)=0$，便于对齐，非理论项。
在 $17$--$31$ 处已见跳跃；未贴合处主要来自有限零点截断。}
\label{fig:riemann-logcos-approx}
\end{figure}

\clearpage

% ============================================================
\section{三角函数与对数余弦：逼近原理上的差异}
% ============================================================

二者都用余弦叠加逼近带跳跃的目标，原理不同。

\begin{enumerate}
  \item \textbf{自变量。}
  $\cos(kx)$ 以 $x$ 为相位，适配加法周期；$L^2$ 上整数谐波是周期函数的自然基。
  $\cos(\gamma\log x)$ 以 $\log x$ 为相位，适配乘法尺度（与 $\pi\sim x/\log x$ 的主趋势一致）。
  前者是通用周期逼近工具；后者来自 $\zeta$ 零点，针对 $\pi-\Li$（或 $\psi-x$）的固有振荡。

  \item \textbf{频率来源。}
  Fourier：先定周期与等距频率，再由目标投影系数——可拟合任意 $L^2$ 样本阶梯。
  黎曼：频率 $\gamma$ 由 $\zeta(\rho)=0$ 给出，振幅含 $x^{\beta}/(|\rho|\log x)$——
  不是对 $\pi(x)$ 的自由拟合。同记「$N$ 项」时，截断的是谐波个数与零点个数，对象不同。

  \item \textbf{误差形态。}
  有限三角和逼近跳跃：经典 Gibbs（过冲比例 $\to$ 常数 $\approx 9\%$）。
  有限零点对数余弦：误差来自未计入的 $\rho$ 及 $\psi\to\pi$ 中略去的项；
  无与 Dirichlet 核相同的那条 Gibbs 定理，但有限 $N$ 下仍可有局部起伏。

  \item \textbf{如何读逼近效果。}
  Fourier 贴阶梯，说明三角系表达力强，不说明已揭示素数—$\zeta$ 机制。
  对数余弦每一项对应零点贡献；$N\to\infty$（对称截断）回到显式公式，并与 RH 下误差阶相连。
  比较「谁更贴」时须分清：通用函数逼近，还是素数显式结构的截断。
\end{enumerate}

简言之：$\cos(kx)$ 服务加法周期；$\cos(\gamma\log x)$ 服务乘法尺度下由 $\zeta$ 规定的素数振荡。
不宜仅用同一 $N$ 的观感判定优劣。

\begin{center}
\renewcommand{\arraystretch}{1.5}
\begin{tabular}{|p{2.8cm}|p{5.2cm}|p{5.5cm}|}
\hline
 & \textbf{Fourier} & \textbf{黎曼对数余弦} \\
\hline
相位 & $x$ & $\log x$ \\
\hline
频率 & 等距整数 & 非等距 $\{\gamma_k\}$ \\
\hline
振幅 & 常系数 & 缓变 $2\sqrt{x}/(|\rho|\log x)$（RH） \\
\hline
截断 $N$ & 谐波个数 & 计入的零点数 \\
\hline
跳跃截断 & Gibbs $\approx 9\%$ & 零点截断起伏（核不同） \\
\hline
\end{tabular}
\end{center}

% ============================================================
\section*{本章小结与后续展望}
\addcontentsline{toc}{section}{本章小结与后续展望}
% ============================================================

共轭零点合并为 $\frac{2x^{\beta}}{|\rho|}\cos(\gamma\log x-\arg\rho)$；
RH 下得 $\psi_0$ 的对数余弦级数~\eqref{eq:cosine_psi}，
$\pi$ 侧主振荡为~\eqref{eq:riemann}（项 $R_k$ 见~\eqref{eq:Rk}）。
$N=3$ 为对数余弦有限截断示意（图~\ref{fig:n3-logcos}）。
Fourier 级数与部分和见~\eqref{eq:fourier}、\eqref{eq:fourier_partial}；
Gibbs 现象见定义~\ref{def:gibbs}（与第~\ref{ch:integral_transforms}~章的 Fourier
\textbf{变换}不同；未展开方波等经典例子）。
$N=100$ 并置对照见图~\ref{fig:fourier-approx}、\ref{fig:riemann-logcos-approx}。
$\cos(kx)$ 与 $\cos(\gamma\log x)$ 分别对应加法周期与乘法尺度；
同记「$N$ 项」时对象不同，误差机制也不同。

\begin{center}
\fbox{\parbox{0.92\textwidth}{
\centering
\vspace{0.2em}
\textbf{第~\ref{ch:fourier_vs_riemann}~章：核心公式}\\[6pt]
\renewcommand{\arraystretch}{1.5}
\begin{tabular}{lp{9.2cm}}
\hline
共轭配对 & \eqref{eq:cosine_pair} \\
RH 下 $\psi_0$ 对数余弦 & \eqref{eq:cosine_psi} \\
$\pi_0-\Li$ 主振荡（近似） & \eqref{eq:riemann} \\
单项修正 $R_k$ & \eqref{eq:Rk} \\
Fourier 级数 / 部分和 & \eqref{eq:fourier}, \eqref{eq:fourier_partial} \\
\hline
\end{tabular}
\vspace{0.2em}
}}
\end{center}

\paragraph{后续展望}
对数余弦截断要真正算出来，须先求临界线上的零点虚部，再代入有限个 $\rho$。
下一章（第~\ref{ch:riemann_numerical_approximation}~章）分两步：Hardy $Z$ 与 Riemann--Siegel 定位 $\gamma$；
显式公式有限截断逼近 $\psi$ 或 $\pi_0$ 并讨论误差。
再后为几何示意（第~\ref{ch:visualization}~章）与 1859 年原稿解读
（第~\ref{ch:riemann_original_interpretation}~章）。
